\section{Coefficients, traces, and division-free identities}
\label{sec:trace}

Throughout this section $R$ is a commutative ring with identity and $A$ is an
$n\times n$ matrix.  All sums and coefficient operations below are finite.
Write
\[
 E(t)=\det(I+tA)=\sum_{k=0}^n e_k t^k,
 \qquad s_j=\operatorname{tr}(A^j)\quad(j\geq1),
 \qquad d_k=(-1)^k e_k.
\]
Set $e_0=d_0=1$, and set $e_k=d_k=0$ for $k>n$.  Then
\[
 p_A(X)=X^n+d_1X^{n-1}+\cdots+d_n.
\]
The coefficient $d_k$ here is the coefficient called $c_{n-k}$ in the
article.  A natural-number multiple such as $k e_k$ means repeated addition;
it does not presuppose that $k$ is invertible.

\subsection{Principal minors and the first coefficients}

\begin{proposition}[Principal-minor formula]
\label{prop:trace-principal}
For $0\leq k\leq n$,
\[
 e_k=\sum_{\substack{S\subseteq\{1,\ldots,n\}\\ |S|=k}}
       \det A[S,S].
\]
In particular $e_1=\operatorname{tr}(A)$, $e_n=\det A$, and
\[
 e_2=\sum_{1\leq i<j\leq n}
       (a_{ii}a_{jj}-a_{ij}a_{ji}).
\]
\end{proposition}
\begin{proof}
Expand the determinant by its permutation definition, and in each factor
$\delta_{i,\sigma(i)}+t a_{i,\sigma(i)}$ choose either summand.  A choice of
the second summand on precisely the rows in $S$ contributes a term of degree
$|S|$.  The other factors vanish unless $\sigma$ fixes every index outside
$S$.  Such permutations are exactly the extensions, by fixed points, of
permutations of the increasingly ordered set $S$.  Their signs are the
signs of those restricted permutations: adjoining a fixed point does not
change the parity of a permutation.  Thus their sum is the determinant of
the principal submatrix indexed by $S$.  This gives the asserted coefficient
by a finite reindexing.  The special cases are its instances for subsets of
size one, two, and $n$.
\end{proof}

\subsection{Newton identities without division or eigenvalues}

\begin{proposition}[Finite Jacobi identity]
\label{prop:trace-jacobi}
For a matrix $M(t)$ of polynomials,
\[
 \frac{d}{dt}\det M(t)
 =\operatorname{tr}\bigl(\operatorname{adj} M(t)\,M'(t)\bigr).
\]
The derivative is the formal derivative of a finite coefficient list.
\end{proposition}
\begin{proof}
Apply the product rule to every product in the permutation formula for the
determinant.  Group terms according to the position $(i,j)$ of the
differentiated factor.  The sum of the remaining signed products is the
cofactor $C_{ij}(M)$, so the result is
$\sum_{i,j}C_{ij}(M)M'_{ij}$.  Since
$(\operatorname{adj}M)_{ji}=C_{ij}(M)$, this sum is the stated trace.
Both the product rule and this regrouping are identities between finite
sums, valid in every characteristic.
\end{proof}

\begin{theorem}[Division-free Newton identities]
\label{thm:trace-newton}
For every $k\geq1$,
\begin{equation}
 k e_k=\sum_{i=1}^{k}(-1)^{i-1}e_{k-i}s_i,
 \qquad
 k d_k+\sum_{i=1}^{k}d_{k-i}s_i=0.
 \label{eq:trace-newton}
\end{equation}
These formulas include the cases $k>n$, using the zero convention above.
\end{theorem}
\begin{proof}
Fix $k$.  The finite geometric sum
\[
 U(t)=\sum_{j=0}^{k-1}(-t)^jA^j
\]
satisfies $(I+tA)U(t)=I+(-1)^{k-1}t^kA^k$.  Consequently, modulo $t^k$,
the adjugate identity gives
\[
 \operatorname{adj}(I+tA)=E(t)U(t).
\]
For completeness, multiply
$\operatorname{adj}(I+tA)(I+tA)=E(t)I$ on the right by $U(t)$ and discard
only coefficients of degree at least $k$; this gives the displayed
congruence without cancelling $E(t)$.  The finite Jacobi identity now yields
\[
 E'(t)\equiv E(t)\sum_{j=0}^{k-1}(-1)^j s_{j+1}t^j\pmod{t^k}.
\]
Comparing the coefficient of $t^{k-1}$ proves the first equality.  Multiplying
it by $(-1)^k$ proves the second.  Differentiation was performed on ordinary
polynomials before truncation; no derivative on the quotient ring
$R[t]/(t^k)$ has been assumed.
\end{proof}

This proof does not require a field, characteristic zero, an eigenvalue, an
algebraic closure, or an inverse of $A$.  It also proves the useful adjugate
trace identities
\begin{equation}
 p_A'(X)=\operatorname{tr}\operatorname{adj}(XI-A),\qquad
 Xp_A'(X)-n p_A(X)
   =\operatorname{tr}\bigl(A\operatorname{adj}(XI-A)\bigr).
 \label{eq:trace-adjugate}
\end{equation}
Indeed the first is Jacobi's formula with derivative matrix $I$; the second
is the trace of
$(XI-A)\operatorname{adj}(XI-A)=p_A(X)I$.

\begin{proposition}[A supplied splitting is enough]
\label{prop:trace-splitting}
Suppose a list $\lambda_1,\ldots,\lambda_n$ in a commutative extension ring
is supplied with the identity
$p_A(X)=\prod_{j=1}^n(X-\lambda_j)$.  Then $e_k$ is the $k$th elementary
symmetric expression in that list and, for every $r\geq1$,
\[
 \operatorname{tr}(A^r)=\sum_{j=1}^n\lambda_j^r.
\]
No diagonalization hypothesis is necessary.
\end{proposition}
\begin{proof}
Coefficient comparison gives the assertion about $e_k$.  Put
$u_r=\sum_j\lambda_j^r$.  Differentiate the finite product
$E(t)=\prod_j(1+t\lambda_j)$ and multiply each factor formally by its
truncated geometric inverse.  The same coefficient argument as above gives
$k d_k+\sum_{i=1}^k d_{k-i}u_i=0$.  In this recurrence the coefficient of
$u_k$ is $d_0=1$.  Subtract the recurrence for $s_k$.  Induction on $k$,
using equality of all preceding $u_i$ and $s_i$, gives $u_k=s_k$.
This last induction needs no division by $k$.
\end{proof}

\subsection{Bell-polynomial formulas with all denominators accounted for}

Define the complete exponential Bell polynomials by the finite recurrence
\begin{equation}
 B_0=1,\qquad
 B_k(x_1,\ldots,x_k)=
  \sum_{i=1}^{k}\binom{k-1}{i-1}x_iB_{k-i}(x_1,\ldots,x_{k-i}).
 \label{eq:trace-bell-rec}
\end{equation}
This recurrence is a definition over $\mathbb Z$ and involves no fractions.
For example, $B_1=x_1$, $B_2=x_1^2+x_2$, and
$B_3=x_1^3+3x_1x_2+x_3$.

\begin{theorem}[Integral Bell identity]
\label{thm:trace-bell}
For every $k\geq0$ in every commutative ring,
\begin{equation}
 k!e_k=B_k\bigl(s_1,-s_2,2!s_3,\ldots,
                    (-1)^{k-1}(k-1)!s_k\bigr).
 \label{eq:trace-bell}
\end{equation}
Thus, if $k!$ is a unit, $d_k$ is $(-1)^k/k!$ times the displayed Bell
polynomial.  In particular the article's determinant formula is obtained
at $k=n$ when $n!$ is a unit.
\end{theorem}
\begin{proof}
Use induction, retaining all previously established identities.  Substitute
$x_i=(-1)^{i-1}(i-1)!s_i$ into the Bell recurrence and substitute
$B_{k-i}=(k-i)!e_{k-i}$ from the earlier stages.  The elementary integer
identity
\[
 \binom{k-1}{i-1}(i-1)!(k-i)!=(k-1)!
\]
reduces the result to
$(k-1)!\sum_{i=1}^k(-1)^{i-1}s_i e_{k-i}$, which equals $k!e_k$ by
\eqref{eq:trace-newton}.  There is no cancellation of $(k-1)!$, so the proof
works in positive characteristic and in rings with additive torsion.
\end{proof}

The first four identities are
\begin{align}
 e_1&=s_1,\notag\\
 2e_2&=s_1^2-s_2,\notag\\
 6e_3&=s_1^3-3s_1s_2+2s_3,\label{eq:trace-lowdim}\\
 24e_4&=s_1^4-6s_1^2s_2+3s_2^2+8s_1s_3-6s_4.\notag
\end{align}
They follow by explicitly evaluating the first four Bell recurrences, so
the displayed formulas supply, respectively, the determinant in dimensions
$1,2,3,4$ after the indicated exact division.  They also give the first four
nonleading characteristic coefficients in all larger dimensions.  For
example, if $2$ is invertible, the dimension-three identity is
\[
 A^3-s_1A^2+\frac{s_1^2-s_2}{2}A-\det(A)I=0.
\]
If $2$ and $3$ are invertible, its dimension-four counterpart is
\[
 A^4-s_1A^3+\frac{s_1^2-s_2}{2}A^2
 -\frac{s_1^3-3s_1s_2+2s_3}{6}A+\det(A)I=0.
\]
Over an arbitrary ring, retain the coefficients $e_k$ from principal minors;
one may multiply the displayed equations by the denominators to obtain valid
consequences, but those multiplied equations need not imply Cayley--Hamilton.

\paragraph{Why the characteristic restriction matters.}
For a prime $p$, the two $p\times p$ matrices $0$ and $I$ over
$\mathbb F_p$ have equal traces of every positive power, namely zero, while
their determinants are respectively $0$ and $1$.  Their characteristic
polynomials are $X^p$ and $X^p-1$.  Thus the traces alone cannot recover the
characteristic coefficients in arbitrary characteristic.  The factor
$k!$ in \eqref{eq:trace-bell} cannot simply be cancelled.

\subsection{The Hessenberg determinant formula}

Let $H_k$ be the $k\times k$ matrix with entries, in one-based indexing,
\[
 (H_k)_{ij}=
 \begin{cases}
 s_{i-j+1},&j\leq i,\\
 k-i,&j=i+1,\\
 0,&j>i+1.
 \end{cases}
\]
Thus its first row is $(s_1,k-1,0,\ldots,0)$ and its last row is
$(s_k,s_{k-1},\ldots,s_1)$.  Put $D_0=1$ and $D_k=\det H_k$.

\begin{theorem}[Integral Hessenberg identity]
\label{thm:trace-hessenberg}
For every $k\geq0$, $D_k=k!e_k$.  In particular, when $k!$ is a unit,
\[
 c_{n-k}=d_k=\frac{(-1)^k}{k!}\det H_k\qquad(0\leq k\leq n).
\]
\end{theorem}
\begin{proof}
Expand $D_k$ along its first column.  In the minor corresponding to row $i$,
the first $i-1$ rows force the entries
$k-1,k-2,\ldots,k-i+1$ successively: the first such row has only its first
available entry nonzero, and the same holds after each removal.  The
remaining lower-right matrix is $H_{k-i}$.  Therefore
\[
 D_k=\sum_{i=1}^k(-1)^{i-1}s_i
       \left(\prod_{h=1}^{i-1}(k-h)\right)D_{k-i}.
\]
The empty product is $1$.  Inductively substituting
$D_{k-i}=(k-i)!e_{k-i}$ turns the product of the two factorial factors into
$(k-1)!$.  Newton's identity now gives $D_k=k!e_k$, as required.
\end{proof}

\subsection{A correctly initialized Faddeev--LeVerrier algorithm}

Assume for now that $1,2,\ldots,n$ are units of $R$.  Define
\begin{equation}
 F_0=I,\qquad
 d_k=-k^{-1}\operatorname{tr}(AF_{k-1}),\qquad
 F_k=AF_{k-1}+d_k I\quad(1\leq k\leq n).
 \label{eq:trace-flv}
\end{equation}
These assignments compute the characteristic coefficients.  Indeed the
invariant
\[
 F_k=A^k+d_1A^{k-1}+\cdots+d_kI
\]
is immediate by induction.  Assuming the coefficients through $k-1$ have
been computed correctly, Newton's identity gives
\[
 \operatorname{tr}(AF_{k-1})
  =s_k+d_1s_{k-1}+\cdots+d_{k-1}s_1=-k d_k,
\]
which proves the next assignment.  The invariant at $k=n$ gives
$F_n=p_A(A)=0$.  The adjugate coefficient formula additionally gives
$\operatorname{adj}(A)=(-1)^{n-1}F_{n-1}$ for $n\geq1$; if $d_n$ is a unit,
$A^{-1}=-d_n^{-1}F_{n-1}$.

For integer matrices the same algorithm uses exact integer division by
the positive number $k$.  Its divisibility is not guessed: the principal-minor
definition supplies the integer witness $-d_k$ for
$\operatorname{tr}(AF_{k-1})=k(-d_k)$.  Integer quotient and remainder are
primitive recursive, and their defining equations verify each step.  More
generally an effective exact-division operation with this correctness
property suffices; mere characteristic zero does not supply such an operation.

The article's displayed recurrence with initial value $M_0=0$ contains
$1/(k-1)$ at $k=1$.  Read literally it cannot be initialized.  One must give
$M_1=I$ and start that recurrence at $k=2$, or use
\eqref{eq:trace-flv} from $F_0=I$.

\subsection{Finite formal logarithms, exponentials, and the partition formula}

Fix a precision $N$ and suppose $1,2,\ldots,N$ are units, as they are in a
$\mathbb Q$-algebra.  Every expression in this subsection is reduced modulo
$z^{N+1}$.  Define
\[
 L_N(z)=-\sum_{j=1}^N\frac{s_j}{j}z^j,
 \qquad
 G_N(z)=\sum_{r=0}^N\frac{L_N(z)^r}{r!}\pmod{z^{N+1}}.
\]
These are finite polynomial definitions, not limits.  Then
\begin{equation}
 \det(I-zA)\equiv G_N(z)
   =\exp_N\bigl(\operatorname{tr}\log_N(I-zA)\bigr)
       \pmod{z^{N+1}}.
 \label{eq:trace-explog}
\end{equation}
To prove the identity, write $g_k$ for the coefficient of $z^k$ in $G_N$.
Differentiating the ordinary finite sum defining $G_N$ gives
$G_N'\equiv L_N'G_N\pmod{z^N}$: the omitted term is
$L_N'L_N^N/N!$ and has degree at least $N$.  Coefficient comparison therefore
gives
\[
 g_0=1,\qquad k g_k=-\sum_{j=1}^k s_j g_{k-j}\quad(1\leq k\leq N).
\]
The coefficients $d_k$ of $\det(I-zA)$ satisfy exactly this recurrence by
Newton's identity.  Induction and multiplication by the supplied inverse of
$k$ give $g_k=d_k$ through degree $N$.  This proves
\eqref{eq:trace-explog} without analytic convergence or a logarithm branch.
No nonzero-eigenvalue assumption is needed.  If $N>n$, the coefficients of
degrees $n+1,\ldots,N$ in $G_N$ consequently vanish.

Expanding the finite exponential and grouping its products according to
the number $m_j$ of occurrences of $s_jz^j/j$ gives the useful corrected
partition formula
\begin{equation}
 d_k=
 \sum_{\substack{m_1,\ldots,m_k\geq0\\
                 \sum_{j=1}^k j m_j=k}}
 \prod_{j=1}^k
   \frac{(-1)^{m_j}s_j^{m_j}}{j^{m_j}m_j!}.
 \label{eq:trace-partition}
\end{equation}
For a fixed tuple put $r=\sum_jm_j$.  Its terms in $L_N^r$ occur
$r!/\prod_jm_j!$ times; division by $r!$ proves the displayed coefficient.
Every index has the bound $m_j\leq k$, so this is a finite bounded sum.
The article's product with sign $(-1)^{m_j+1}$ for each of $n$ factors differs
from the correct sign by $(-1)^n$; for odd $n$ it would already give the
incorrect leading coefficient $c_n=-1$.  Equation
\eqref{eq:trace-partition} uses the corrected sign.  In rings where the
fractions are unavailable, the integral Bell identity is the applicable
statement instead.

The familiar expression
$p_A(\lambda)=\lambda^n\exp(\operatorname{tr}\log(I-A/\lambda))$
can therefore be read rigorously, at every finite precision, by putting
$z=\lambda^{-1}$ in \eqref{eq:trace-explog}.  This is a formal Laurent
calculation.  It does not assert the existence of a numerical matrix
logarithm for every nonzero value of $\lambda$.

\subsection{The exact induction obligations for iPRA}

For the arithmetic implementation, subsets are bit masks below $2^n$,
permutations are bounded lists, and polynomials are finite coefficient
arrays with an explicit length.  Traces, determinants, products, powers,
coefficient selection, formal derivatives, and truncation are primitive
recursive operations on these codes.  The earlier ring-identity
interpretation is used for the universal ring assertions.  Rational
instances use signed integer numerator--positive denominator pairs;
statements about supplied inverses include their finite verification
equations.

Computability by itself is not the proof of iPRA-provability.  The preceding
proofs use the following explicit induction invariants, each a decidable
equality of finite arrays or a bounded conjunction of such equalities.
For a coded ring presentation use its decidable equality; for generic
identities compare the corresponding normalized integer polynomial arrays.
\begin{enumerate}
\item The geometric-sum invariant is
$(I+tA)\sum_{j=0}^{q}(-t)^jA^j
 =I+(-1)^qt^{q+1}A^{q+1}$.
Its successor step is obtained by adjoining one term and cancelling the
two adjacent terms of degree $q+1$.
\item Determinant differentiation is a bounded sum of product-rule
identities, and the principal-minor argument is the explicit finite
bijection $(\sigma,S)\leftrightarrow(S,\sigma|_S)$ on nonzero summands.
Finite-sum reindexing is proved by induction on the enumerated list length;
zero summands are retained, so no existence search for a nonzero entry occurs.
\item In the Bell proof retain the table invariant
$B_j=j!e_j$ for every $j\leq q$.  Its next row follows by
\eqref{eq:trace-bell-rec} and \eqref{eq:trace-newton}.  For the Hessenberg
proof retain $D_j=j!e_j$ on the same prefix.  Their factorial identities are
equalities of finite products and binomial coefficients, with the usual
Pascal recursion giving quantifier-free induction proofs.
\item In the Faddeev--LeVerrier proof retain the polynomial expression for
$F_q$ and correctness of the first $q$ coefficients.  The next numerator
is exactly $-k d_k$ by Newton's identity, so unit multiplication or verified
exact integer division proves the next step.
\item In the finite exponential proof retain $g_j=d_j$ for $j\leq q$.
The displayed coefficient recurrence and the explicit inverse of $q+1$
give the successor step.  The partition expansion is a bounded grouping of
words of length at most $N$ by their occurrence-count vectors.
\end{enumerate}
A prefix invariant can be represented by one primitive recursive Boolean
conjunction over its bounded list.  Induction on that Boolean equation is
quantifier-free induction, so no induction over an unbounded quantified
formula is being silently introduced.  These invariants give the concrete
proof obligations for the object-language formalization, rather than
inferring provability merely from the existence of the algorithms.
